\begin{abstract}
LLM Agents can hide which tool operation a user requested by sending decoy calls. When an observer sees the hop to a third-party tool and the privacy target includes which request--response exchanges occur and how long the service takes, TLS and HTTP/2 padding cannot imitate responses, so decoys must reach the application. We show that such cover leaks through the process that dispatches it. Honest tool calls follow a model decision; decoys inserted by a controller do not wait for one, and this difference is exploitable on real TLS records. In 3,000 single-host executions with two fixed-weight planners over MCP and TLS~1.3/HTTP/2, label-free timing removal raises recognition of the requested action from 75 to 112 (Qwen) and from 80 to 108 (Llama) of 120 tasks when call boundaries are supplied. Dispatching decoys after sampled historical gaps defeats that removal, yet adapted attacks still gain over size-only attacks in three of four comparisons after Holm correction, recognizing 93--106/120. A shadow planner on an independent ticket brings the strongest tested attack to 70/120 and 64/120, consistent with a 7/12 plausible-deniability reference, at twice the model tokens and with more and different cover calls. We define an evaluation that binds inference, task completion, business-state integrity and cost to the same executions, and state which service capabilities each policy assumes.
\end{abstract}
